\section{Discussion and Conclusion}
\label{sec:conclusion}
We have developed new bounds based on the sigmoid function that allow tractable gradient-based optimization of any differentiable classifier to meet a desired precision constraint. Our experiments evaluating in-hospital mortality prediction on two large hospital datasets suggest the feasibility of our approach in real early warning applications.

\textbf{Limitations.} A key drawback of our method is vulnerability to local optima due to non-convexity of the sigmoid bound.
Given reasonable compute power at training time, we find we can mitigate local optima by taking the best of many runs.
%For MLPs, cold starts using Glorot initialization~\citep{glorot2010understanding} satisfy the false alarm contraint more often than warm starts from a BCE solution (see Tables~\ref{table:mimic_performance}-\ref{table:eicu_performance}).

% The chosen value of $\alpha$, the desired minimum precision to enforce, is critical. While the $\alpha$ values used in our experiments are reasonable, in practice $\alpha$ needs to be chosen in collaboration with stakeholders. 

For a deployment, setting $\alpha$ requires input from stakeholders (at what false alarm rate would alerts be ignored or not add value?). 
As in most real-world ML, we expect a modest but noticeable drop between train and heldout performance.
As a rule of thumb, if the desired heldout precision is 0.5, we recommend setting $\alpha$ a bit higher (say $\alpha = 0.6$) for training then enforce $\alpha' = 0.5$ on a validation set to be sure desired performance is achieved.
Our method is sensitive to other hyperparameters, including the penalty $\lambda > 0$ and the tolerances $\gamma, \delta, \epsilon$ that govern the tractability-tightness tradeoff for our bound. 
While we found reasonable values for our tasks, future applications may need to tune values for better performance.

%In this work, we fixed the parameters of the sigmoid bound to ensure a reasonably tight bound, and also ran our experiments with a pre-defined $\alpha$. In practice however, $\alpha$ needs to be chosen in collaboration with experienced clinical staff, and $\lambda, \gamma, \delta, \epsilon$ need to be searched over a hyperparameter grid to guarantee optimal performance. The vast space of hyper-parameters would require large computational capacity. 

Finally, although our objective can control false alarms on the training set, this does not guarantee this constraint will be satisfied on heldout data.
Guaranteeing generalization is an open research question. 
%Perhaps  tightening the bounds gradually during training could overcome this issue.


\textbf{Advantages.}
The primary advantage of our method is the ability to steer models to satisfy a maximum desired false alarm rate, in order to avoid alarm fatigue, while maintaining linear runtime.
As demonstrated throughout our experiments, our objective clearly outperforms alternatives like post-hoc threshold search at maximizing recall at a desired minimum precision.

% Our bounds can work with any classifier trained via SGD. 
Future work could use our bounds on more flexible classifiers that are trainable via SGD, such as recurrent NNs, convolutional NNs, or graph NNs.
For improved interpretability, one could explore integrating our loss with L0 and L1 norm weight sparsity penalties~\citep{tibshirani1996regression, ustunSupersparseLinearInteger2016}.
